%============================================================ SECTION 1: INTRODUCTION
\section{Introduction \& Recap of the Design}

\subsection{The attack and its objective}
Network reconnaissance is the opening move of any intrusion: before an attacker can act,
they must learn which services a target exposes, what software versions run there, and
which operating system underpins them. Project~8 asks us to build this capability
(port scanning with OS information and version) entirely from scratch, crafting every frame,
packet, and segment in our own code.

Our toolkit carries out a full reconnaissance-to-exploitation chain in five phases:
(1)~raw TCP~SYN port scanning, (2)~service banner grabbing, (3)~probabilistic OS
fingerprinting, (4)~OS-aware payload selection, and (5)~a dynamic exploit battery,
paired with (6)~a defensive detector that recognizes the scan and can quarantine the
attacker in real time.

\subsection{What changed since the Design Report}
The week-11 Design Report proposed a functional but smaller system. During
implementation the project expanded across three distinct areas.
This report is organized around that progression: each core section compares the
originally planned behavior with the delivered implementation. The table
below summarizes the additions.

\renewcommand{\arraystretch}{1.12}
\noindent\begin{tabularx}{\textwidth}{@{}L{3.2cm} L{5.2cm} X@{}}
\toprule
\hdrrow \thh{Area} & \thh{Design Report (week 11)} & \thh{Delivered implementation}\\
\midrule
OS fingerprinting &
Single passive SYN/ACK feature table (TTL, window, MSS, WScale, SACK, timestamp, option order). &
\textbf{Seven-strategy probabilistic engine}: RST analysis, IP-ID sequencing,
timestamp-clock recovery, ECN probe, invalid-flag probe, ICMP/UDP battery, cross-layer
banner correlation, with confidence scoring and coarse version tiers.\\
\addlinespace[3pt]
Exploit suite &
Three fixed proof-of-concept exploits (FTP anon, SSH, HTTP). &
\textbf{Five OS-aware modules} with dynamic payload selection: FTP anonymous, FTP
write$\rightarrow$web-shell$\rightarrow$RCE, FTP bounce, HTTP traversal, and SSH
brute-force with quiet/loud post-exploitation.\\
\addlinespace[3pt]
Active defense &
Single fast sliding-window SYN detector; Windows-only \texttt{netsh} block. &
\textbf{Four-vector, cross-platform engine}: fast + slow windows, invalid-flag detection,
canary honeyport; \texttt{netsh} \emph{and} \texttt{iptables}; self-healing TTL release,
whitelist, and JSON incident logging.\\
\addlinespace[3pt]
Test topology &
One generic Windows victim. &
\textbf{Five real operating systems}: Windows~11, Kali~Linux, Android, iPadOS/Darwin,
embedded-Linux router.\\
\addlinespace[3pt]
Verification &
Manual observation. &
\textbf{68 automated tests} plus twelve archived empirical run logs.\\
\bottomrule
\end{tabularx}
\renewcommand{\arraystretch}{1.0}

\subsection{Lab topology}
\noindent
\begin{minipage}[t]{0.45\textwidth}
\vspace{0pt}
The verification testbed is a single isolated private LAN. The attacker is a Windows~11 laptop running the toolkit over an Npcap adapter; the primary victim is a Windows~11 desktop hosting the vulnerable lab services; and a set of secondary targets on the same subnet exercises the fingerprinting engine against genuinely different TCP/IP stacks.

\vspace{3pt}
The attacker never leaves this private range: an RFC~1918 check refuses any non-private target, and the scanner is bounded (Section~\ref{sec:safety}) so that it cannot be repurposed as a flooding tool.
\end{minipage}\hfill
\begin{minipage}[t]{0.53\textwidth}
\vspace{0pt}
\centering
\resizebox{\linewidth}{!}{%
\begin{tikzpicture}[
  font=\sffamily\scriptsize, >={Stealth[length=1.6mm]},
  box/.style={rectangle, rounded corners=2pt, draw=rule, line width=0.6pt,
              fill=codebg, minimum height=0.68cm, align=center, inner sep=2.5pt},
  atk/.style={box, draw=accent, line width=1.0pt, fill=accentsoft},
  vic/.style={box, draw=accent, line width=1.0pt},
  hub/.style={rectangle, rounded corners=2pt, draw=accent, line width=0.8pt,
              fill=white, align=center, inner sep=3pt, minimum height=0.75cm}]

% attacker (left)
\node[atk] (atk) at (0,0) {\textbf{Attacker}\\[0.5pt]\tiny Win 11 Laptop\\[0.5pt]\tiny scanner $\cdot$ fp $\cdot$ exploits};

% LAN hub (centre)
\node[hub] (sw) at (2.9,0) {\textbf{Isolated LAN}\\[0.5pt]\tiny 192.168.x.0/24};

% targets (right column, compact)
\node[vic]  (vic)  at (6.2,1.35)  {\textbf{Victim}: Win 11\\[0.5pt]\tiny FTP $\cdot$ HTTP $\cdot$ SSH $\cdot$ Def};
\node[box] (kali) at (6.2,0.45)  {Kali VM \tiny(Netfilter)};
\node[box] (mob)  at (6.2,-0.45) {Android $\cdot$ iPadOS};
\node[box] (rtr)  at (6.2,-1.35) {Embedded Router};

% attacker -> hub
\draw[->,line width=0.8pt,accent] (atk) -- node[above=1pt,font=\tiny\color{accent}]{raw SYN} (sw);

% hub -> each target
\draw[->,line width=0.6pt,rule] (sw) -- (vic);
\draw[->,line width=0.6pt,rule] (sw) -- (kali);
\draw[->,line width=0.6pt,rule] (sw) -- (mob);
\draw[->,line width=0.6pt,rule] (sw) -- (rtr);
\end{tikzpicture}%
}
\vspace{-2pt}
\captionof{figure}{Heterogeneous verification testbed. One attacker drives probes across five distinct OS stacks on an isolated LAN.}
\label{fig:topology}
\end{minipage}
